% [DATA-SOURCE] Rec-01..Rec-10 | TN3K routing_log.jsonl de-identified exemplars
\noindent\textbf{去标识化路由记录样例（Rec-01--Rec-10）。}
\label{sec_A1_routing_exemplars}
\par\smallskip
\noindent\textbf{Rec-01.}\par
\begin{itemize}[leftmargin=1.5em,topsep=0.1em,itemsep=0.05em]
\item 病例标识: Rec-01
\item 任务类型: 分类路由（分类器不一致时的 LLM 裁决）
\item 候选相近度分层: 相近
\item 是否调用 LLM: 是
\item 决策方式: LLM 裁决
\item 最终良恶性类别: 良性
\item 置信度: 0.720
\item 决策说明: 独立分类输入：4 个独立分类模型: 恶性 (0.66)、良性 (0.78)、良性 (0.99)、良性 (1.00)；AutoGluon（所选 mask）：恶性（0.587）
\end{itemize}
\noindent\textbf{Rec-02.}\par
\begin{itemize}[leftmargin=1.5em,topsep=0.1em,itemsep=0.05em]
\item 病例标识: Rec-02
\item 任务类型: 分割路由
\item 候选相近度分层: 相近
\item 是否调用 LLM: 是
\item 决策方式: LLM 选 mask
\item 所选分割专家: TN3K 专用
\item 置信度: 0.840
\item 决策说明: 候选：多队列联合（一致度 0.820）；TN3K 专用［所选］（一致度 0.819）；ThyroidXL 专用（一致度 0.733）；PKTN 专用（一致度 0.772）；TN5K 专用（一致度 0.664）；分类锚点：良性
\end{itemize}
\noindent\textbf{Rec-03.}\par
\begin{itemize}[leftmargin=1.5em,topsep=0.1em,itemsep=0.05em]
\item 病例标识: Rec-03
\item 任务类型: 分类路由（分类器不一致时的 LLM 裁决）
\item 候选相近度分层: 相近
\item 是否调用 LLM: 是
\item 决策方式: LLM 裁决
\item 最终良恶性类别: 恶性
\item 置信度: 0.650
\item 决策说明: 独立分类输入：4 个独立分类模型: 恶性 (0.64)、恶性 (1.00)、良性 (0.96)、良性 (1.00)；AutoGluon（所选 mask）：恶性（0.766）
\end{itemize}
\noindent\textbf{Rec-04.}\par
\begin{itemize}[leftmargin=1.5em,topsep=0.1em,itemsep=0.05em]
\item 病例标识: Rec-04
\item 任务类型: 分割路由
\item 候选相近度分层: 相近
\item 是否调用 LLM: 是
\item 决策方式: LLM 选 mask
\item 所选分割专家: TN3K 专用
\item 置信度: 0.680
\item 决策说明: 候选：多队列联合（一致度 0.355）；TN3K 专用［所选］（一致度 0.369）；ThyroidXL 专用（一致度 0.317）；PKTN 专用（一致度 0.301）；TN5K 专用（一致度 0.090）
\end{itemize}
\noindent\textbf{Rec-05.}\par
\begin{itemize}[leftmargin=1.5em,topsep=0.1em,itemsep=0.05em]
\item 病例标识: Rec-05
\item 任务类型: 分类路由（分类器不一致时的 LLM 裁决）
\item 候选相近度分层: 相近
\item 是否调用 LLM: 是
\item 决策方式: LLM 裁决
\item 最终良恶性类别: 恶性
\item 置信度: 0.600
\item 决策说明: 独立分类输入：4 个独立分类模型: 恶性 (1.00)、恶性 (1.00)、良性 (0.81)、良性 (1.00)；AutoGluon（所选 mask）：恶性（0.661）
\end{itemize}
\noindent\textbf{Rec-06.}\par
\begin{itemize}[leftmargin=1.5em,topsep=0.1em,itemsep=0.05em]
\item 病例标识: Rec-06
\item 任务类型: 分割路由
\item 候选相近度分层: 相近
\item 是否调用 LLM: 是
\item 决策方式: LLM 选 mask
\item 所选分割专家: PKTN 专用
\item 置信度: 0.860
\item 决策说明: 候选：多队列联合（一致度 0.886）；TN3K 专用（一致度 0.916）；ThyroidXL 专用（一致度 0.893）；PKTN 专用［所选］（一致度 0.917）；TN5K 专用（一致度 0.856）；分类锚点：恶性
\end{itemize}
\noindent\textbf{\textbf{Rec-07}.}\par
\begin{itemize}[leftmargin=1.5em,topsep=0.1em,itemsep=0.05em]
\item 病例标识: Rec-07
\item 任务类型: 分割路由
\item 候选相近度分层: 相近
\item 是否调用 LLM: 是
\item 决策方式: LLM 选 mask
\item 所选分割专家: 多队列联合
\item 置信度: 0.860
\item 决策说明: 候选：多队列联合［所选］（一致度 0.795）；TN3K 专用（一致度 0.796）；ThyroidXL 专用（一致度 0.776）；PKTN 专用（一致度 0.803）；TN5K 专用（一致度 0.499）；分类锚点：良性
\end{itemize}
\noindent\textbf{Rec-08.}\par
\begin{itemize}[leftmargin=1.5em,topsep=0.1em,itemsep=0.05em]
\item 病例标识: Rec-08
\item 任务类型: 分类路由（分类器不一致时的 LLM 裁决）
\item 候选相近度分层: 相近
\item 是否调用 LLM: 是
\item 决策方式: LLM 裁决
\item 最终良恶性类别: 恶性
\item 置信度: 0.600
\item 决策说明: 独立分类输入：4 个独立分类模型: 良性 (0.97)、恶性 (0.83)、良性 (0.93)、良性 (1.00)；AutoGluon（所选 mask）：恶性（0.768）
\end{itemize}
\noindent\textbf{Rec-09.}\par
\begin{itemize}[leftmargin=1.5em,topsep=0.1em,itemsep=0.05em]
\item 病例标识: Rec-09
\item 任务类型: 分类路由（分类器不一致时的 LLM 裁决）
\item 候选相近度分层: 相近
\item 是否调用 LLM: 是
\item 决策方式: LLM 裁决
\item 最终良恶性类别: 恶性
\item 置信度: 0.830
\item 决策说明: 独立分类输入：4 个独立分类模型: 恶性 (1.00)、恶性 (1.00)、恶性 (0.93)、良性 (1.00)；AutoGluon（所选 mask）：恶性（0.779）
\end{itemize}
\noindent\textbf{Rec-10.}\par
\begin{itemize}[leftmargin=1.5em,topsep=0.1em,itemsep=0.05em]
\item 病例标识: Rec-10
\item 任务类型: 分类路由（确定性共识）
\item 候选相近度分层: 易区分
\item 是否调用 LLM: 否
\item 决策方式: 确定性共识
\item 最终良恶性类别: 良性（共识）
\item 置信度: 0.998
\item 决策说明: 独立分类模型一致预测为「良性」(conf=0.95\textasciitilde{}1.00)，直接采纳。
\end{itemize}
\par\smallskip\noindent\footnotesize 以上十条去标识化样例（Rec-01--Rec-10）在分割 LLM 路由与分类路由（含 LLM 调用与确定性共识等）之间均衡选取；所列分割 LLM 路由的相近度分层均为「相近」（最优与次优候选一致度差 $\leq 0.10$）。样例字段与表~A1-2 一致。
